\scanpage{123}
\begin{flushright}\fbox{\shortstack{Complexes de \uline{\unsure{chaînes|chaines|duaux}} de De Rham\\à p.d.}}\end{flushright}
Soient

$k$ anneau comm

\deleted{$\unsure{S|A}$} \textcircled{$s$} une $k$-coalgèbre \inserted{ass. unitaire comm}

\deleted{\unsure{avec $k$}}
\[
 \boxed{s\longrightarrow s\otimes_k s}
\]
[donc \inserted{$V_k(\deleted{\unsure{S}}s):$} $k'\longmapsto s^{k'}=\operatorname{Hom}_k(s,k')$ est un foncteur à valeurs dans les anneaux ass. unit. comm, \unsure{soit|fait} une $\mathcal O_k$-Alg. ass. un. comm

\diagram{\begin{tikzcd}[ampersand replacement=\&,column sep=small]
 |[draw,circle,inner sep=1pt]|{\bar s}\arrow[d,shift left]\arrow[from=2-1,to=1-1,shift left] \& \text{une sous-}k\text{-coalgèbre \inserted{ass. unit. comm}}\\
 |[draw,circle,inner sep=1pt]|{s^+} \& {}=s/\bar s\quad\text{le }k\text{-module quotient}
\end{tikzcd}}
[donc \inserted{$V_k(\bar s):$} $k'\longmapsto\bar s^{k'}=\operatorname{Hom}_k(\bar s,k)$ est une $\mathcal O_k$-alg. ass. unit. comm et $s^{k'}\longrightarrow\bar s^{k'}$ est un hom de $k'$-Alg. ass. unit. comm \inserted{i.e. $V_k(s)\to V_k(\bar s)$ hom de $\mathcal O_k$-algs.}

[c'est une coalgèbre ass. comm mais pas unitaire]

[$s^{+\,k'}=\operatorname{Ker}(s^{k'}\to\bar s^{k'})\subset s^{k'}$ est un idéal de $s^{k'}$, donc $V_k(s^+)\overset{\text{déf}}=V_k(s)^+$ est un Idéal de $V_k(s)$]

\textcircled{$\gamma_*$} un système de copuissances divisées sur $(s,s^+)$

[donc on a sur $V_k(s^+)=V_k(s)^+$ un syst. de puiss. divisées
\[
 \gamma_i:V_k(s^+)\longrightarrow V_k(s^+),\quad i\geq1
\]
défini par des hom de $k$-alg
\[
 \gamma_i^*:\operatorname{Sym}_k(s^+)\longrightarrow\operatorname{Sym}_k^*(s^+)
\]
\deleted{\unsure{qui}} i.e. hom $k$-\unsure{lin|lin.}
\[
 \gamma_i^*:s^+\longrightarrow\operatorname{Sym}_k^*(s^+)
\]
qui, en vertu des axiomes sont en fait
\[
 \boxed{\gamma_i^*:s^+\longrightarrow\operatorname{Sym}_k^i(s^+)\quad(i\geq1)}
\]
avec des axiomes \unsure{dessus|divers}

$t:s^+\to k$ un $k$ hom \qquad i.e. $t\in V_k(s^+)(k)=V_k(s)^+(k)$

Alors on définit un objet simplicial dans la \unsure{cat.} des foncteurs
\[
 (k\text{-alg comm})\longrightarrow
 \left(\begin{aligned}
 &\text{algèbres graduées}\\[-2pt]
 &\text{alternées différentielles}\\
 &\text{(diff. de degré }+1)\text{ à}\\
 &\inserted{\text{avec idéal à}}\text{ puiss. divisées }\unsure{(|)}
 \end{aligned}\right.
\]
et $=$ dans la \unsure{cat.} des \deleted{foncteurs} schémas affines \deleted{$\unsure{S}$} en algèbres graduées alternées différentielles (diff. de degré $+1$) avec Idéal \inserted{homogène} à \unsure{puiss.|puissances divisées} (qui est un schéma affine), les puiss. divisées étant compatibles aux graduations et à la différentielle
% archive sheet number: 122
